\section{Limitations}\label{sec:limits}

\paragraph{The theory.} Model~F is an idealization: it assumes that every turn misses, that requests complete in FIFO order and
that a finished turn rejoins at once. The engine runs about 28 requests concurrently, vLLM preempts, and hits return as the backlog
falls toward the onset, so model~F describes the trap, not the transitions into and out of it. We do not prove that the engine
converges to model~F, and the theorem's limits are reached only geometrically: in the registered window about 33--37 generations
after the burst, a transient of order 10\% remains, so we compare with the exact fluid as well as with the limits. The premium is
proved positive for our workload at $\rho\ge1.04$; below that it rests on a floating-point bracket and a first-order expansion.

\paragraph{The validation.} The check on the 57 stuck runs is retrospective: the theory was derived after they were read, and its
scoring tolerances were set on them. Only the runs submitted after tag \texttt{icml-trap-pred} test it blind
(\S\ref{sec:validation}). The from-empty runs are single runs per rate, the trapped windows are at most 7500\,s long, and ``does not
escape'' means within that horizon.

\paragraph{The experiments.} Sessions are synthetic (random tokens, exponential 5\,s tool times, geometric turn counts, fixed tool
outputs and output lengths, no shared prefixes), clients never retry or abandon, and real agent traces have tool times far shorter
and more dispersed. Engines are colocated, with one GPU each, dense 7--8B models and pools of 150k to 956k tokens; disaggregated
serving, offloading, larger models and parallelism are untested. Seeds are shared across tests. The rescue engages a fixed cap at
a fixed time, and the arrival stop and the tool slowdown ran at one rate. Appendix~\ref{sec:limits-old} gives the full list.

\section{Conclusion}\label{sec:conclusion}

A burst can tip a prefix-cached agent server into a trap in which every turn misses; Lu et al.\ and we observed this
independently. We have shown that the trap has an exact law. Idealized as a fluid FIFO server on which every turn misses, the
trapped engine obeys a linear renewal recursion along FIFO orbits, so from any trapped state it converges to a self-similar regime
fixed by a Perron root: the gap between turns grows linearly, the backlog grows at a fixed fraction of the excess load, the turns it
serves are young and short, and it completes sessions faster than its all-miss capacity. Measured in 57 stuck runs on two engines,
two GPUs and two models, with nothing fitted to trap data, these laws hold. The practical reading is Corollary~\ref{cor:trapthr}:
the throughput an operator sees in a trap overstates the rate to which arrivals must be cut, and cutting below the all-miss
capacity drains the trap only slowly and by turning sessions away. The other exit, a session cap that restores reuse, rescued every
established trap we tried at the same load without shedding, and the engine stayed good after the cap was removed.
